% OWNER: sections/s2_3_queries.tex. Edit only this file.
\seclabel
\claim{\secnum{2.3}{sec:queries}\quad A non-trading date returns another day's data with HTTP 200}

\lead{Every REST request has the form \code{/api/{area}/{name}}, and the paths come from code served only after login.} The public JavaScript bundle names only the \code{Auth} paths. The trading screen's code lives under \code{/assets/app/}, which is served only with the \code{asset_token} cookie (\figref{creds}), so we took the other paths from the code downloaded after login (\figref{paths}). The administrator areas (\code{Master}, \code{Admin}) are outside the wrapper's scope.

% OWNER: figs/paths.tex. Edit only this file.
\fig{paths}{The wrapper's REST paths fall in 5 areas; bold paths were called}[Main REST paths extracted from the web client's code, grouped by their first path segment]{%
{\tabfont\begin{tabularx}{\linewidth}{>{\RaggedRight\arraybackslash}p{16mm} >{\RaggedRight\arraybackslash}X}
\thd{Area} & \thd{Paths (selection)}\\
\code{Auth} & \textbf{\code{Login}}, \textbf{\code{SilentRefresh}}, \textbf{\code{Logout}}\\\rowsep
\code{Portfolio} & Orders: \textbf{\code{AddOrder}}, \code{ModifyOrder}, \code{CancelOrder}, \code{DeleteOrder}, \textbf{\code{CancelOrderBatch}}, \textbf{\code{ValidateProduct}}, \code{OrderLogs}, \textbf{\code{Orders}}\\
 & State: \textbf{\code{PfBs}}, \textbf{\code{Session}}, \textbf{\code{Trades}}, \code{Fills}, \code{FillAggrs}, \textbf{\code{Universe}}, \textbf{\code{Violations}}, \code{ViolationsSummary}, \code{PerfCharts}\\\rowsep
\code{Contest} & \code{List}, \code{PfList}, \code{TopRankDetail}, \code{TopPick}\\\rowsep
\code{Analysis} & \code{PfAsOf}, \code{PrftByProd}, \code{HistPosChart}, \code{TrendChart}, \code{ActivityChart}\\\rowsep
\code{Util} & \code{SearchProducts}\\
\end{tabularx}}}{A full path joins area and name, e.g.\ \code{Portfolio/Orders}. Bold: called live on 24 Nov 2026 (\code{PfBs} answered 404, \figref{readeps}); roman: read from the client code only. Request bodies of the order paths are given in \figref{orderapi}.}{Web client code, DART Quant Team}


\lead{A Bearer header is enough; queries need no cookies.} A \code{GET} to \code{Portfolio/Violations} for the live portfolio returned HTTP 200 with the access token in the header and 401 without it.

\lead{The live portfolio has to be chosen by contest ID, reset flag and end date, because 2 portfolios can share a contest name.} The \code{pfs} list gains an entry for each contest entered and each reset (\figref{pfs}). Round 13 has 2 entries with the same contest name, and picking by name can return the 1--8 Oct entry, which ended with a reset.

% OWNER: figs/pfs.tex. Edit only this file.
\fig{pfs}{The live portfolio is neither reset nor past its end date}[The 6 entries in \code{pfs} from the login reply, newest first; portfolio IDs omitted]{%
{\tabfont\begin{tabularx}{\linewidth}{>{\RaggedRight\arraybackslash}p{30mm} >{\RaggedRight\arraybackslash}p{25mm} >{\RaggedRight\arraybackslash}p{13mm} >{\RaggedRight\arraybackslash}X}
\thd{Contest (\code{ctstId})} & \thd{Period} & \thd{\code{invalid}} & \thd{Status on the web page}\\
\hl\textbf{Round 13 (91)} & \textbf{12 Oct to 30 Nov} & \textbf{absent} & \textbf{live}\\\rowsep
Round 13 (91) & 1 Oct to 8 Oct & True & ended, reset\\\rowsep
Round 12 (90) & 3 Jul to 31 Aug & absent & ended\\\rowsep
Round 12 (90) & 1 Jul to 2 Jul & True & ended, reset\\\rowsep
Round 11 (86) & 1 Apr to 29 May & absent & ended\\\rowsep
Practice (1) & 5 Feb to 30 Apr & True & ended, reset\\
\end{tabularx}}}{Shaded row: the live portfolio, the only entry without \code{invalid} whose period includes 24 Nov. \code{invalid} appears only on reset portfolios. \code{active} is False for all 6 entries, including the live one, so the wrapper selects by contest (\code{ctstId}), absent \code{invalid} and an end date (\code{lastD}) not yet passed, and requires exactly 1 match.}{DART Quant Team}


\lead{Two field names mislead.} \code{active} is False for all 6 entries, including the live one, so it says nothing about which portfolio is running. The reset flag is the presence of \code{invalid}, so a missing key has to be read as False. If the rule below matches anything other than exactly one portfolio, the wrapper stops instead of guessing.

\par\noindent\begin{minipage}{\linewidth}\fontsize{8.5}{11}\selectfont\ttfamily
\hspace*{4mm}select\_pf(ctstId):\\
\hspace*{8mm}cands $\leftarrow$ \{ pf $\in$ pfs : pf.ctstId = ctstId $\wedge$ $\neg$pf.invalid\\
\hspace*{38mm}$\wedge$ pf.lastD $\geq$ currBizDay \}\\
\hspace*{8mm}return the single element of cands, else raise\end{minipage}\par

\lead{The date parameter accepts any date, and the server decides which day's data to return (\figref{dates}).} A weekend date returns the previous trading day, a future date the latest one, and a holiday after the contest returns zeros, all with HTTP 200. A caller that collects history day by day therefore has to check each date against a trading calendar, or the same day's data is stored under several dates; the wrapper itself has no such loop.

% OWNER: figs/dates.tex. Edit only this file.
\fig{dates}{Weekend and future dates return another day's data without an error}[\code{Portfolio/Violations} replies by requested date, 2026 (server business day: 23 Nov)]{%
{\tabfont\begin{tabularx}{\linewidth}{>{\RaggedRight\arraybackslash}p{44mm} R{12mm} R{11mm} >{\RaggedRight\setlength{\leftskip}{3mm}\arraybackslash}X}
\thd{Requested date} & \thd{HTTP} & \thd{Bytes} & \thd{Day the server used, and the evidence}\\
22 Nov (Sunday) & 200 & 400 & 20 Nov, the previous trading day: the small-cap weight is 3.21\%, against 3.33\% on 23 Nov\\\rowsep
23 Nov (current day) & 200 & 400 & 23 Nov\\\rowsep
24 Nov (next trading day) & 200 & 400 & 23 Nov, the latest: byte-for-byte identical to the 23 Nov reply\\\rowsep
25 Dec (holiday, after the competition) & 200 & 342 & No day: every value is zero\\
\end{tabularx}}}{HTTP: status code of the reply. Bytes: reply body size. Replies compared with \code{cmp}; where they differed, we printed the differing fields. All 4 requests succeed, so a caller cannot tell from the status which day the data belong to.}{DART Quant Team}


\lead{The Sunday reply carries 20 Nov data: it differs from the 23 Nov reply only in the weight of the portfolio's one small-cap stock.} That stock's price did not change on 23 Nov (previous close and close both KRW~14,840), while the portfolio's value fell 3.59\% that day. A stock whose price stays flat gains weight as the portfolio shrinks, so its 20 Nov weight should be $3.33 \times (1 - 0.0359) = 3.21$\%, the value the Sunday request returned.

\lead{The server's business day also differs from the calendar for several hours after midnight.} \code{currBizDay} still read 23 Nov at 06:02:10 on 24 Nov and read 24 Nov at 07:06:44, so the server moves to the next business day between those two times, before the 09:00 open. A request made before the rollover carries the previous day's date.

\lead{Three more read endpoints give the order window, the investable universe and the order history (\figref{readeps}).} \code{Session} returns 09:00 to 15:20, the window in which orders execute; 15:20 is when the closing auction begins. \code{Universe} lists 2,648 investable stocks with their GICS sector. Some tickers mix letters and digits (\code{A0015M0}), following the Korea Exchange's newer code format, so tickers are handled as strings. \code{Orders} takes a date range, \code{d0} and \code{d1}. Our first request sent \code{d} and got an empty list for 23 Nov, although the live feed showed a filled order on that day; the web page's own request during the session sent \code{d0} and \code{d1} and received every order of the range, finished ones included (Section~\ref{sec:orders}).

% OWNER: figs/readeps.tex. Edit only this file.
\fig[wide]{readeps}{Four read endpoints answer; \code{Orders} needs a date range, \code{PfBs} returns 404}[\code{Portfolio} read endpoints called with the live \code{pfId} and \code{d} (07:33 KST), \code{Orders} with \code{d0} and \code{d1} (09:16 KST), 24 Nov 2026]{%
{\tabfont\begin{tabularx}{\linewidth}{>{\RaggedRight\arraybackslash}p{14mm} R{10mm} R{16mm} >{\RaggedRight\setlength{\leftskip}{3mm}\arraybackslash}X >{\RaggedRight\arraybackslash}p{52mm}}
\thd{Path} & \thd{HTTP} & \thd{Records} & \thd{Content} & \thd{Use in the wrapper}\\
\code{Session} & 200 & 1 & \code{hm0} 09:00, \code{hm1} 15:20 & Bounds for order start and end times\\\rowsep
\code{Universe} & 200 & 2,648 & \code{Id}, \code{nm}, \code{sec} (GICS sector); some tickers alphanumeric & Ticker validation, sector of each stock\\\rowsep
\code{Trades} & 200 & 62 & 1 record per position from entry to exit; 23 still open at 07:33, matching the 23 holdings; 64 records after the day's entries; times in UTC & Position history and realised profit\\\rowsep
\code{Orders} & 200 & 1 per order & Every order from \code{d0} to \code{d1}, \code{Done} ones included, so the count depends on the dates sent. Fields: \code{Id}, \code{state}, \code{qty2ord}, \code{cumQty}, \code{avgPx} (average fill price, KRW), \code{cnclT}, \code{genT}, \code{lastMCV}, \code{sgn}. Sent with \code{d} instead, it returned an empty list & Day's order count at start-up\\\rowsep
\code{PfBs} & 404 & 0 & Not found under the 3 areas tried (\code{Portfolio}, \code{Contest}, \code{Analysis}), as GET or POST & Not used; balances come from the live feed\\
\end{tabularx}}}{GET with query parameters unless noted. HTTP: status code; Records: items in the reply. \code{Orders} parameters copied from the web page's own request.}{DART Quant Team}

\verdict{the wrapper fills every request date from the server's \code{currBizDay} and refuses to place orders while \code{currBizDay} differs from the KST calendar date (D6, D27). It chooses the portfolio with the rule above and stops when the rule matches 0 or 2 portfolios (D7).}
